\documentclass[12pt]{article}
\usepackage[utf8]{inputenc}
\usepackage[T1]{fontenc}
\usepackage{hyperref}
\title{Pinai-Hagahai Language Resources for K-12 Curriculum Development}
\author{Compiled by Claude Sonnet 5.5 (claude-sonnet-5-5) for Andr\'as Kornai}
\date{}
\begin{document}
\maketitle

\section{Introduction}
Pinai-Hagahai (ISO 639-3 pnn) is a Piawi (Madang -- Upper Yuat) language of Papua New Guinea, spoken along the Yuat River in Enga and Madang Provinces. Speaker figures conflict: 600 (1996), about 700 (2005), 1,500 (Joshua Project). Alternate names include Hagahai, Aramo, Pinai. Search found no school materials.

\section{Primary Educational Resources}
None found in search results.

\section{Linguistic and Reference Materials}
\begin{itemize}
\item Hagahai language, Wikipedia: \url{https://en.wikipedia.org/wiki/Hagahai_language}
\item Piawi languages, Wikipedia: \url{https://en.wikipedia.org/wiki/Piawi_languages}
\item Endangered Languages Project entry: \url{https://endangeredlanguages.com/lang/10684?hl=en}
\end{itemize}

\section{Religious and Community Resources}
\begin{itemize}
\item Joshua Project, Pinai in Papua New Guinea: \url{https://joshuaproject.net/people_groups/14395}
\end{itemize}

\section{Digital and Audio Resources}
\begin{itemize}
\item Wikitongues page for pnn (content not verified): \url{https://wikitongues.org/languages/pnn/}
\end{itemize}

\section{Curriculum Development Context}
The language is still widely spoken, so a mother-tongue-based programme is plausible, but search surfaced only reference pages. Orthography and materials would need to be developed with the community.
\end{document}